\documentclass{article}
\usepackage[utf8]{inputenc}
\usepackage{amsmath,amssymb}
\usepackage[most]{tcolorbox}
\usepackage{geometry}
\geometry{margin=2.5cm}

\newcommand{\step}[1]{\par\noindent\hangindent=3.5em\hangafter=1%
  \makebox[3.5em][l]{\textbf{#1.}}}
\newcommand{\steptag}[1]{\hfill{\small\textsf{[#1]}}}

\newtcolorbox{proofbox}[1]{
  colback=gray!5, colframe=gray!60,
  fonttitle=\bfseries, title={#1},
  breakable, enhanced,
  left=4pt, right=4pt, top=4pt, bottom=4pt,
  before skip=10pt, after skip=10pt
}

\begin{document}

\begin{proofbox}{After first fixing one global witness package W\_RF suppli...}
\step{1} After first fixing one global witness package W\_RF supplied by lem-routef-raw-factor-setting-formation, for every input (H,Phi,eta) to which that formation result applies, fix one def-routef-raw-factor-setting datum S over that same W\_RF supplied by the same result, for every Delta' supplied for (W\_RF,S) by lem-routef-delta-prime-closeness, every Delta supplied from that same Delta' by lem-routef-delta-normalization-closeness, every Upsilon' supplied from that same pair by lem-routef-upsilon-prime-closeness, and every Upsilon supplied from that same triple by lem-routef-upsilon-normalization-closeness, writing the fields of (W\_RF,S) as the unqualified symbols below: Route F common coefficient/domain: K in (1.6) is finite and universal, and rho\_fac in (1.7) is positive and is a common domain for the degree-two estimate and the three Route-F factorization estimates. \steptag{validated} \\[6pt]
\step{1.1} The prefix of node 1 is a nested conditional binding, not an assertion that Delta', Delta, Upsilon', and Upsilon exist for every formation-applicable eta. The formation lemma supplies one global W\_RF and, for each formation-applicable input, a datum S; node 1 then quantifies successively over any Delta', Delta, Upsilon', and Upsilon supplied on the respective provider domains, with the displayed dependency relations and unqualified notation. Moreover, on the common domain 0 <= eta <= rho\_fac where the four estimates are invoked, all four providers are applicable in sequence. \steptag{validated} \\[6pt]
\step{1.1.1} The contract prefix literally has the nested form: after fixing W\_RF and S, for every Delta prime supplied by the first provider, every Delta supplied from that Delta prime by the second provider, every Upsilon prime supplied from that same pair by the third provider, and every Upsilon supplied from that same triple by the fourth provider. Therefore it binds any such supplied maps conditionally and makes no unconditional existence claim for eta outside the providers respective radius hypotheses. \steptag{validated} \\[4pt]
\step{1.1.2} If 0 <= eta <= rho\_fac, then (1.7) gives eta <= rho\_2 and eta <= rho\_DeltaUpsilon. From (1.2), rho\_2 <= rho\_Delta prime and rho\_2 <= rho\_Delta, so the Delta-prime and Delta providers apply successively. From (1.5), rho\_DeltaUpsilon <= rho\_Upsilon, and the definition of rho\_Upsilon in (1.5) gives rho\_Upsilon <= rho\_Upsilon prime; hence the Upsilon-prime and Upsilon providers also apply successively. Thus all four maps exist with the required dependency relations on the common domain, while outside that domain node 1 asserts only the nested conditional quantification. \steptag{validated} \\[4pt]
\step{1.2} For the globally fixed header W\_RF, every coefficient entering K is a finite positive real scalar independent of H, Phi, eta, dimension, amplification level, and block data, and each radius entering rho\_fac is positive. \steptag{validated} \\[6pt]
\step{1.2.1} By lem-routef-raw-factor-setting-formation and (1.1) of def-routef-raw-factor-setting, C\_theta, C\_A, bar\_C\_E, C\_V, and C\_T are finite positive real scalars; rho\_theta, rho\_AI, epsilon\_E/C\_A, [4(1+C\_theta)]\textasciicircum{}(-1), and [4(1+C\_V)]\textasciicircum{}(-1) are positive, so their finite minimum rho\_T is positive. All these scalars are independent of the input data because W\_RF is one globally fixed independent header and the displayed operations use only its fields and numerical constants. \steptag{validated} \\[4pt]
\step{1.2.2} Starting from the positive finite scalars of (1.1), the serial definitions in (1.2) make C\_Delta', C\_Delta, C\_2, and C\_3 finite positive, and make rho\_unit, rho\_id, rho\_prod, rho\_Delta', rho\_Delta, rho\_2, rho\_DeltaPhi, and rho\_3 positive: sums and positive integer multiples preserve finite positivity, reciprocals have strictly positive denominators, and every displayed minimum has finitely many positive entries. These derived scalars remain independent of all input data. \steptag{validated} \\[6pt]
\step{1.2.2.1} Directly expanding (1.1)-(1.2) of def-routef-raw-factor-setting and using lem-routef-raw-factor-setting-formation: C\_theta=12*(sqrt(2)-1)>0, C\_A=20+(211/8)*C\_theta>0, bar\_C\_E=max\{1,C\_E\}>=1, hence C\_V=bar\_C\_E*C\_A>0 and C\_T=C\_theta+3*C\_V>0, all finite; eta\_A, epsilon\_E, rho\_theta, and rho\_AI are positive, so every entry of rho\_T is positive. Substitution into (1.2) then proves, in displayed order, the asserted coefficient finiteness/positivity and radius positivity, because each operation is a finite sum, positive multiple, reciprocal of a positive scalar, or finite minimum of positive scalars. Global-header independence is preserved under these scalar operations. \steptag{validated} \\[4pt]
\step{1.2.3} Using the positive finite quantities already produced by (1.1)-(1.2), (1.3) makes C\_N, C\_R, C\_L, and C\_Upsilon' finite positive; since C\_R>0, (2*C\_R)\textasciicircum{}(-1)>0, and therefore (1.4) makes rho\_Upsilon' a finite minimum of positive radii and hence positive. These derived scalars remain independent of all input data. \steptag{validated} \\[6pt]
\step{1.2.3.1} Direct expansion from the formation data through (1.1)-(1.3) of def-routef-raw-factor-setting gives finite positive C\_V, C\_T, C\_Delta', C\_Delta, C\_2, and C\_3; therefore the four sums defining C\_N, C\_R, C\_L, and C\_Upsilon' are finite positive. Direct expansion of the radius formulas through (1.2) gives positive rho\_T, rho\_id, rho\_Delta, rho\_2, and rho\_3, while C\_R>0 gives (2*C\_R)\textasciicircum{}(-1)>0; their minimum rho\_Upsilon' in (1.4) is thus positive. Every value uses only the fixed input-independent header and numerical scalar operations. \steptag{validated} \\[4pt]
\step{1.2.4} Using the positive finite quantities produced by (1.1)-(1.4), (1.5) makes C\_Upsilon finite positive and makes rho\_Upsilon, rho\_DeltaUpsilon, rho\_mult, and rho\_UpsilonDelta positive: both new reciprocal denominators are strictly positive and every displayed minimum has only positive entries. These derived scalars remain independent of H, Phi, eta, dimension, amplification level, and block data. \steptag{validated} \\[6pt]
\step{1.2.4.1} Direct expansion from lem-routef-raw-factor-setting-formation through (1.1)-(1.4) of def-routef-raw-factor-setting yields finite positive C\_T and C\_Upsilon' and positive rho\_unit, rho\_Upsilon', rho\_theta, rho\_T, rho\_id, rho\_Delta, and rho\_DeltaPhi. Hence C\_Upsilon=6*C\_T+7*C\_Upsilon' is finite positive, [2*(C\_T+C\_Upsilon')]\textasciicircum{}(-1)>0, and each minimum defining rho\_Upsilon, rho\_DeltaUpsilon, rho\_mult, and rho\_UpsilonDelta in (1.5) is positive. These quantities depend only on the globally fixed input-independent header and numerical scalar operations. \steptag{validated} \\[4pt]
\step{1.3} The scalar K defined by (1.6) is finite and universal. \steptag{validated} \\[6pt]
\step{1.3.1} By lem-routef-raw-factor-setting-formation, the fields of W\_RF are one fixed package independent of H, Phi, eta, dimension, amplification level, and block data. Direct serial expansion of (1.1)-(1.5) in def-routef-raw-factor-setting uses only finite real sums, products, maxima, minima, and reciprocals with positive denominators, so C\_theta, C\_Delta, C\_Upsilon, and C\_2 are finite and input-independent. Consequently all four entries in the maximum (1.6) are finite and input-independent, and their finite maximum K is finite and universal. \steptag{validated} \\[4pt]
\step{1.4} The scalar rho\_fac defined by (1.7) is positive. \steptag{validated} \\[6pt]
\step{1.4.1} From lem-routef-raw-factor-setting-formation, eta\_A>0 and epsilon\_E>0, while C\_theta=12*(sqrt(2)-1)>0, C\_A=20+(211/8)*C\_theta>0, rho\_theta=1/8>0, and bar\_C\_E=max\{1,C\_E\}>=1. Expanding (1.1)-(1.5) of def-routef-raw-factor-setting in displayed order, every coefficient used as a reciprocal denominator is positive and every radius is a finite minimum of positive quantities; in particular rho\_2, rho\_DeltaUpsilon, rho\_mult, and rho\_UpsilonDelta are positive. Equation (1.7) defines rho\_fac as the minimum of exactly these four positive numbers, hence rho\_fac>0. \steptag{validated} \\[4pt]
\step{1.5} Every eta with 0 <= eta <= rho\_fac lies simultaneously in the domains of lem-routef-degree-two-estimate, lem-routef-delta-upsilon-telescope, lem-routef-multiplicative-telescope, and lem-routef-upsilon-delta-telescope; hence rho\_fac is a common domain for the degree-two estimate and the three Route-F factorization estimates. \steptag{validated} \\[6pt]
\step{1.5.1} Equation (1.7) of def-routef-raw-factor-setting gives rho\_fac<=rho\_2, rho\_fac<=rho\_DeltaUpsilon, rho\_fac<=rho\_mult, and rho\_fac<=rho\_UpsilonDelta. Thus 0<=eta<=rho\_fac implies each corresponding radius hypothesis. In the ambient package supplied respectively by lem-routef-raw-factor-setting-formation, lem-routef-delta-prime-closeness, lem-routef-delta-normalization-closeness, lem-routef-upsilon-prime-closeness, and lem-routef-upsilon-normalization-closeness, exact application of lem-routef-degree-two-estimate at rho\_2, lem-routef-delta-upsilon-telescope at rho\_DeltaUpsilon, lem-routef-multiplicative-telescope at rho\_mult, and lem-routef-upsilon-delta-telescope at rho\_UpsilonDelta yields all four estimates on 0<=eta<=rho\_fac. Therefore rho\_fac is their common domain. \steptag{validated} \\[4pt]
\end{proofbox}

\end{document}
